\documentclass{subfiles}
\begin{document}
    We now turn our attention to a more interesting question: 
        fixed a source over some alphabet, how do we design the optimal code for it?
        The question, apparently hard, has a somewhat obvious answer:
        take in account the lengths of the codewords. Or better said, 
        the average one.

    Formally speaking, we take in account the \gls{acl} of the code, 
        defined as the quantity
        \[
            L_{S}(C) = \sum_{i = 1}^{n} p_{i}l_{i}
        \]
        such that is minimal. We say that a code \(\C*\) is compact (or optimal) for the source 
        \(S\) if, among all prefix codes for \(S\), \(\C*\) has the minimal ACL.

    \begin{remark*}
        Is there a lower bound for the ACL? Yes, and it given by the following result.
        \begin{lemma*}[Shannon]
            If \(\C*\) is a UD code for a source \(S\) with probabilities 
            \(p_{1}, \ldots, p_{n}\), and \(d\) is the size of the code alphabet,
            then
            \[
                L_{S}(\C*) \ge \frac{H(S)}{\log d}
            \]
        \end{lemma*}
    \end{remark*}

    In addition to the ACL, we can define two addtional quantities that gives informations 
        on the code: the efficiency \(\vartheta(\C*) = \tfrac{H(S)}{L_{S}(\C*)}\),
        which by Shannon lemma is bounded between 0 and 1; 
        and the redundancy \(\rho_{S}(\C*) = 1 - \vartheta_{S}(\C*)\).

    Some attempt to provide a way to build optimal codes have been tried:
        the Shannon-Fano encoding for instance; but none achived the goal.
        \clearpage
\end{document}
